\section{Global Gauss area and the cost of marked collars}
\label{sec:ed03-global-gauss-caps}

\begin{theorem}[Oriented normal and the ordinary area of the direction sphere]
\label{thm:ed03-global-normal}
\uses{lem:ed03-normal-heights}
\lean{SmoothPoincare.orientedGaussVector_volume_pos, SmoothPoincare.orientedGaussNormal_smooth, SmoothPoincare.orientedPrincipalFrameAt_middle_order, SmoothPoincare.gaussRoundArea_eq_gaussConeArea, SmoothPoincare.gaussRoundArea_univ, SmoothPoincare.gaussRoundArea_antipodal}
\leanok
Let $F:(M^4,g)\to\mathbb R^5$ be a smooth oriented isometric hypersurface
immersion. Fix the ambient orientation by an orthonormal basis. The unique
positive unit normal defines a smooth map $N:M\to S^4$. At each point
there is an oriented orthonormal principal frame for this normal with
nondecreasing signed principal curvatures. Ordinary Riemannian area on
the Euclidean unit sphere agrees with radial cone area, has total mass
$\omega=8\pi^2/3$, and is invariant under the antipodal map.
\end{theorem}
\begin{proof}
\leanok
Normalize the oriented cross product of the four tangent vectors of an
oriented chart. A positive change of basis leaves its normalization fixed,
also where the second fundamental form vanishes. Diagonalize the symmetric
shape operator, order its eigenvectors, and adjust the source orientation.
For the area calculation use $(u,t)\mapsto e^tu$. The pulled-back ambient
volume has density $e^{5t}$ relative to round area times $dt$, with integral
$1/5$ over $t<0$. Comparing the volumes of cones gives equality of the
two measures on every Borel set.
\end{proof}

\begin{theorem}[Open middle locus and the indexed area identity]
\label{thm:ed03-ordinary-indexed-area}
\uses{thm:ed03-global-normal,prop:ed03-gauss-sheets}
\lean{SmoothPoincare.isOpen_middlePrincipalLocus, SmoothPoincare.gaussParamDensity_eq_mul_frameJacobian, SmoothPoincare.ImmersionGaussNormal.measuredSheets_of_principal, SmoothPoincare.orientedImmersion_indexedGaussMean}
\leanok
On a nonempty compact Polish source modeled on Euclidean four-space, let $U$
be the locus of signature $(2,2)$ and $J=|\prod_{j=1}^4 k_j|$. Write
$\mu$ for ordinary round area divided by $\omega$. For a unit direction $u$
let $m_2(u)$ count the nondegenerate index-two critical points of
$\langle F,u\rangle$ on the source. The locus $U$ is open, $J$ is
continuous, the point set is finite almost surely, and $m_2$ is measurable
and integrable, with
\[
    \int_{S^4}m_2(u)\,d\mu(u)=\frac{2}{\omega}\int_U J\,dV_g .
\]
\end{theorem}
\begin{proof}
\leanok
In a single smooth normal chart use the determinant of the signed second
form, and two fixed tangent coordinate vectors where the quadratic values
have opposite strict signs. Positive determinant excludes degeneracy and
forces an even number of negative eigenvalues. The two evaluations exclude
both definite signatures. All three strict tests persist near the center.
Changing the unit normal preserves middle signature. In the same source
parametrization the target density of the composed Gauss map equals the
source density times its orthonormal absolute Jacobian. Ratios on charts
give continuity of $J$ even at singular points. Integrate on countably many
regular sheets and use the disjoint fibers over $u$ and $-u$. Multiplicity
counts points in $M$, including distinct points with the same immersed image.
\end{proof}

\begin{theorem}[Morse directions at the immersed threshold]
\label{thm:ed03-morse-threshold}
\uses{thm:ed03-ordinary-indexed-area,cor:ed03-strict-actual-height}
\lean{SmoothPoincare.riemannianVolumeMeasure_criticalValues_eq_zero, SmoothPoincare.orientedHeight_ae_nondegenerate, SmoothPoincare.inducedEnergy_strict_of_orientedHeightFloor, SmoothPoincare.exists_morseHeight_le_one_of_inducedEnergy_le}
\leanok
Almost every round direction gives a smooth height with nondegenerate
Hessian at each critical point. On the compact source above, use the
ordered frames of the constructed normal. With integrable squared
chiral norm and continuous middle deficit, energy at most $64\pi^2/9$
gives a Morse height with a finite set of at most one index-two point.
\end{theorem}
\begin{proof}
\leanok
Apply the equal-dimensional Sard theorem in charts and transfer the null
critical values to Riemannian volume. Remove both this set and its antipode.
The signed Weingarten formula then gives a nondegenerate Hessian at
every critical point of the height. If almost every direction had at
least two middle points, the indexed identity and exclusion of open
balanced principal frames would force strict energy above the threshold.
Intersect the conull Morse and finiteness conclusions before selecting
a direction. Coincident critical values are allowed.
\end{proof}
Smooth recognition from one such direction still requires a marked handle
presentation and its boundary gluings on the original smooth source.

\begin{proposition}[Physical chart tests and moving frames]
\label{prop:ed03-chart-frame-tests}
\uses{prop:ed03-anchored-shell,lem:ed03-splice-curvature}
\lean{SmoothPoincare.physicalRadialVolume_eq_map_scaleMetric, SmoothPoincare.physicalAnchor_and_anchoredBuffer_fullCurvature_tendsto_zero, SmoothPoincare.shiftedPhysicalEuclideanShellVolume_strict_decrease, SmoothPoincare.exists_orientedMovingFrame, SmoothPoincare.movingCurvatureSlots_tensor, SmoothPoincare.tendstoUniformlyOn_movingChiralProjector}
\leanok
On a supplied wide chart, radial approximation on the closed buffered
band gives a physical anchor and transfers the curvature ceiling to
that band. Already in Euclidean space a positive shift of the radial
coordinate gives a strictly decreasing normalized adjacent-shell volume
quotient. On one oriented coordinate basis, positive metric matrices
admit a continuous choice of compatible orthonormal frames. On a fixed
uniform ellipticity window uniform convergence of the metric coefficients
gives uniform convergence of the density-weighted chiral projections.
\end{proposition}
\begin{proof}
\leanok
Push forward the restricted physical volume by the normalized physical
distance; retain the omitted region when comparing whole shells.
For the shifted quotient apply the mean value theorem to differences of
fourth powers. On the compact matrix window apply uniform continuity
to orthonormalization and the four tensor slots in this same basis.
\end{proof}

\begin{proposition}[Positive marked caps and their individual costs]
\label{prop:ed03-cap-tests}
\uses{prop:ed03-inclusive-caps,lem:ed03-splice-curvature}
\lean{SmoothPoincare.HasJoinedCapTests.yamabeFloor, SmoothPoincare.twoCapSigns_of_joinedTests, SmoothPoincare.HasJoinedCapTests.markedSide, SmoothPoincare.cappedEnergy_le_raw_add_halfCollarError}
\leanok
Set $T=4\pi^2$. On one connected joined path, assume the continuous sign
test, positive start, Euler bound and energy ceiling $2T-\gamma$ for
$\gamma>0$. If every smooth log-conformal test on each completed cap
has joined transplants with both numerator and volume converging,
each cap has Yamabe constant at least $\sqrt{48\gamma}$. The actual
cap metric, raw side and complementary marked core then define a positive
marked side. On that cap, strip mass $s$, norm error $\beta$, and a comparison
of the entire output complement give a cap cost at most
$E_{\mathrm{raw}}+2\sqrt{s}\,\beta+\beta^2$.
\end{proposition}
\begin{proof}
\leanok
Pass the joined lower bound to each test quotient before taking the
infimum. The limiting test volume is positive. For the cost, apply the
$L^2$ triangle inequality on the trim and add the complementary integral.
The two caps use their own half-collar errors. Once both signs and markings
are available, the sum test in Proposition~\ref{prop:ed03-inclusive-caps}
includes equality.
\end{proof}

\begin{proposition}[The selected disk widths and finite band charge]
\label{prop:ed03-same-width}
\uses{prop:ed03-endpoint-windows,prop:ed03-cap-tests}
\lean{SmoothPoincare.endpointPhysical_omitted_charges, SmoothPoincare.exists_selectedDiskWidths_small_chiral_sum, SmoothPoincare.endpointAllowance_on_selectedChartDisks, SmoothPoincare.bandPhysical_omitted_charges, SmoothPoincare.bandPhysical_yamabeEuler_sandwich}
\leanok
Use the volume-one physical Yamabe sequence with the ordered spectral,
Euler and signature data and the modified probability gauge with its
constant test and same-region spectral transport. At $I=T$ and
$e_i\to T$, suppose marked disjoint chart disks capture all but vanishing
modified volume. Their full physical exterior has chiral charge $r_i\to0$.
There is one $k_i\ge i$ per finite family for which the sum of physical
trim charges tends to zero. Write $s_i$ for their union charge at width
$1/(k_i+1)$. For any nonnegative allowance $\eta_i$ on these widths,
suppose that eventually the strict joined ceiling
$e_i+2\sqrt{s_i}\eta_i+\eta_i^2<2T$ yields a positive competitor on
the original type costing at most $(\sqrt{r_i}+\eta_i)^2$. Then for each
$\delta>0$, eventually $\eta_i>2\pi-\delta$.

At a single index, if $U$ is a supplied retained set with omitted modified
probability at most $d$, and $\ell^2$ is the total modified spectral second
moment with $\ell\ge0$, then each physical chiral charge satisfies
\[
 E_\pm(g;U^c)\le 2e-\tfrac32\max\{Y/6-\ell\sqrt d,0\}^2 .
\]
\end{proposition}
\begin{proof}
\leanok
The endpoint scalar test and spectral deficit transfer vanishing omitted
volume to vanishing physical charge. The marked coordinate spheres have
zero physical volume. Thin their collars together at each finite index
and apply the repair test at those widths. A fixed saving below $2\pi$
would violate the fixed-type infimum. At a single index Cauchy--Schwarz
bounds the omitted first moment by $\ell\sqrt d$. Transport the captured
second moment on $U$ and use the common total chiral energy of either
sign. The auxiliary positive-Ricci metric supplies the disk selection in
the manuscript; the charges here are measured in $g$.
\end{proof}
